\section{Chapter 3: Laplace transforms \texorpdfstring{\fr{86--134}}{}}
\label{ch3}

\subsection*{Big picture}
The Laplace transform turns \textbf{differentiation into multiplication by $s$} and \textbf{convolution into multiplication}.
So ODEs and integral equations become algebra: transform, solve for $\hat y(s)$, use partial fractions, read the answer from the table
(justified by Lerch's uniqueness theorem). For $(A,B,C,D)$ it gives $\hat Y=T(s)\hat U$.

\subsection{Definition and the table \texorpdfstring{\fr{86--90}}{}}\label{ss:3-1}
\begin{plainbox}
The Laplace transform is a \textbf{translator}: it turns a function of time $f(t)$ into a function of $s$, $\hat f(s)$.
The point is that hard calculus in $t$ becomes easy algebra in $s$. You will mostly use the \textbf{table}, not the integral.
\par\smallskip\textbf{In the exam this looks like:} ``Using the table, write down $\Lap(\dots)$'' or ``invert $F(s)$''.
\end{plainbox}
\begin{defbox}
$f:(0,\infty)\to\CC$ continuous is of \textbf{exponential type} (satisfies (E)) if $|f(x)|\le Me^{\beta x}$ for some $M>0$, $\beta\in\RR$.
Then
\[\boxed{\Lap(f)(s)=\hat f(s)=\int_0^\infty f(x)e^{-sx}\,dx\qquad(\Re s>\beta).}\]
\end{defbox}
\begin{center}
\renewcommand{\arraystretch}{1.55}
\begin{tabular}{@{}ll|ll@{}}
\toprule
$f(t)$ & $\hat f(s)$ & $f(t)$ & $\hat f(s)$\\\midrule
$1$ & $\dfrac1s$ & $\sin at$ & $\dfrac{a}{s^2+a^2}$\\
$t^n$ & $\dfrac{n!}{s^{n+1}}$ & $\cos at$ & $\dfrac{s}{s^2+a^2}$\\
$e^{at}$ & $\dfrac1{s-a}$ $(s>a)$ & $t^ne^{\lambda t}$ & $\dfrac{n!}{(s-\lambda)^{n+1}}$\\
$t^\alpha$ $(\alpha>-1)$ & $\dfrac{\Gamma(\alpha+1)}{s^{\alpha+1}}$ & $e^{at}\cos bt$ & $\dfrac{s-a}{(s-a)^2+b^2}$\\
$H(t-b)$ & $\dfrac{e^{-bs}}s$ & $e^{at}\sin bt$ & $\dfrac{b}{(s-a)^2+b^2}$\\
$\delta_b$ & $e^{-bs}$ & $\exp(tA)$ & $(sI-A)^{-1}$\\
\bottomrule
\end{tabular}
\end{center}
(The last three of the right-hand column follow from the shift rule and Prop 0.34.)

\subsection{Properties \texorpdfstring{\fr{91--98}}{}}\label{ss:3-2}
\begin{plainbox}
The properties are the ``grammar'' of the translator. The most important: \textbf{differentiating in $t$ becomes multiplying by $s$}
(minus the starting value). So an ODE becomes an equation you can solve by dividing.
\par\smallskip\textbf{In the exam this looks like:} ``State and prove $\Lap(f')=s\Lap f-f(0)$'' or using the shift and $t f(t)$ rules to find a transform.
\end{plainbox}
\begin{thmbox}[title={Proposition 0.25 and 0.26}]
If $f$ satisfies (E) with constants $M,\beta$:
\begin{enumerate}
\item $\Lap f(s)$ exists for $s>\beta$ and $|\Lap f(s)|\le\dfrac{M}{s-\beta}$; hence $\Lap f(s)\to0$ as $s\to\infty$.
\item \textbf{Shift:} $\Lap\big(e^{ax}f(x)\big)(s)=\Lap f(s-a)$.
\item \textbf{Linearity:} $\Lap(\lambda f+\mu g)=\lambda\Lap f+\mu\Lap g$.
\item \textbf{Multiplication by $x$:} $\Lap\big(xf(x)\big)(s)=-\dfrac{d}{ds}\Lap f(s)$.
\item \textbf{Derivative:} $\Lap(f')(s)=s\Lap f(s)-f(0)$; repeatedly,
$\Lap(f'')=s^2\hat f-sf(0)-f'(0)$, \ $\Lap(f''')=s^3\hat f-s^2f(0)-sf'(0)-f''(0)$.
\item $\Lap f$ is \textbf{holomorphic} on $\{\Re s>\beta\}$.
\end{enumerate}
\end{thmbox}

\begin{proofbox}[title={Proof of the derivative rule (5) \fr{93}}]
Integrate by parts on $[0,R]$ for $s>\beta$:
$\int_0^Re^{-sx}f'(x)dx=\big[e^{-sx}f(x)\big]_0^R+s\int_0^Re^{-sx}f(x)dx=e^{-sR}f(R)-f(0)+s\int_0^Re^{-sx}f(x)dx$.
Since $|f(R)|\le M'e^{\beta R}$ (f satisfies (E) because $f(x)=f(0)+\int_0^xf'$), $e^{-sR}f(R)\to0$. Let $R\to\infty$: $\Lap(f')=s\Lap f-f(0)$.
\end{proofbox}
\begin{proofbox}[title={Proof of the bound (1) and the shift rule (2)}]
(1) $\big|\int_0^\infty e^{-sx}f(x)dx\big|\le\int_0^\infty Me^{(\beta-s)x}dx=\dfrac{M}{s-\beta}$ for $s>\beta$.\\
(2) $\Lap(e^{ax}f)(s)=\int_0^\infty e^{-sx}e^{ax}f(x)dx=\int_0^\infty e^{-(s-a)x}f(x)dx=\Lap f(s-a)$.
\end{proofbox}

\subsection{Laplace transforms from the definition}\label{sec:lapdef}\label{ss:3-3}
\begin{plainbox}
``From the definition'' means you must actually do the integral $\int_0^\infty f(t)e^{-st}dt$, not quote the table.
Two tricks cover almost everything: \textbf{integration by parts} (for powers of $t$) and \textbf{writing $\cos,\sin$ using $e^{\pm i a t}$}.
\par\smallskip\textbf{In the exam this looks like:} ``Calculate $\Lap(\cos2\omega t)$ from the definition, deduce $\Lap(\sin^2\omega t)$'' (AE3, A3.1).
\end{plainbox}
\begin{methodbox}
\begin{enumerate}
\item Write $\int_0^R f(x)e^{-sx}dx$ (finite $R$ first, as in the notes).
\item For powers $x^n$: integrate by parts $n$ times. For $\cos,\sin$: write them using $e^{\pm iax}$, or integrate by parts twice and solve for the integral.
\item Let $R\to\infty$, stating the condition on $s$ that makes the boundary terms vanish (e.g.\ $s>0$).
\end{enumerate}
\end{methodbox}

\begin{workbox}[title={Worked example \fr{88}: $\Lap(x^2)=2/s^3$}]
$\int_0^Rx^2e^{-sx}dx=\Big[\frac{x^2e^{-sx}}{-s}\Big]_0^R+\frac2s\int_0^Rxe^{-sx}dx
=\Big[\frac{x^2e^{-sx}}{-s}\Big]_0^R+\Big[\frac{2xe^{-sx}}{-s^2}\Big]_0^R+\Big[\frac{2e^{-sx}}{-s^3}\Big]_0^R
=-\frac{R^2e^{-Rs}}{s}-\frac{2Re^{-Rs}}{s^2}-\frac{2e^{-Rs}}{s^3}+\frac2{s^3}\to\frac2{s^3}$ for $s>0$.
\end{workbox}

\begin{workbox}[title={Worked example (AE3, A3.1): $\Lap(\cos2\omega t)$ and $\Lap(\sin^2\omega t)$}]
\textbf{(i)} Write $\cos2\omega t=\frac12(e^{2i\omega t}+e^{-2i\omega t})$. For $s>0$,
$\int_0^\infty e^{\pm2i\omega t}e^{-st}dt=\Big[\frac{e^{-(s\mp2i\omega)t}}{-(s\mp2i\omega)}\Big]_0^\infty=\frac{1}{s\mp2i\omega}$
(the boundary term at $\infty$ vanishes since $|e^{-(s\mp2i\omega)t}|=e^{-st}\to0$). Hence
\[\Lap(\cos2\omega t)=\frac12\Big(\frac1{s-2i\omega}+\frac1{s+2i\omega}\Big)=\frac12\cdot\frac{2s}{s^2+4\omega^2}=\boxed{\frac{s}{s^2+4\omega^2}}\quad(s>0).\]
\textbf{(ii)} $\sin^2\omega t=\frac12(1-\cos2\omega t)$, so by linearity
\[\Lap(\sin^2\omega t)=\frac1{2s}-\frac{s}{2(s^2+4\omega^2)}=\frac{(s^2+4\omega^2)-s^2}{2s(s^2+4\omega^2)}=\boxed{\frac{2\omega^2}{s(s^2+4\omega^2)}}.\]
\end{workbox}

\begin{workbox}[title={Worked example: Heaviside function \fr{90} and the $t$-rule}]
$\Lap(H(x-b))=\int_b^\infty e^{-sx}dx=\frac{e^{-sb}}{s}$ $(s>0)$.\\
$\Lap(t\sin at)=-\frac{d}{ds}\frac{a}{s^2+a^2}=\frac{2as}{(s^2+a^2)^2}$. \quad $\Lap(t^2e^{-3t})=\frac{2}{(s+3)^3}$ (shift rule applied to $2/s^3$).
\end{workbox}

\subsection{Partial fractions and inverting \texorpdfstring{\fr{129--134}}{}}\label{ss:3-4}
\begin{plainbox}
To go back from $s$ to $t$, break the fraction into \textbf{simple pieces} (partial fractions), because each simple piece is in the table.
Lerch's theorem guarantees the answer is unique, so ``looking it up'' is a valid method.
\par\smallskip\textbf{In the exam this looks like:} the last step of every Laplace question: ``use partial fractions at the final step''.
\end{plainbox}
\begin{thmbox}[title={Prop 0.35, Cor 0.36, Cor 0.37}]
Every rational $f(s)=q(s)+\sum_j\frac{a_j}{(s-\lambda_j)^{n_j}}$ uniquely (Euclidean algorithm then partial fractions).
$\Lap\big(t^ne^{\lambda t}\big)=\dfrac{n!}{(s-\lambda)^{n+1}}$, so \textbf{strictly proper rational functions are exactly the Laplace transforms of sums of $t^ne^{\lambda t}$}.
The Laplace transform is a bijection between $\mathrm{span}\{t^ne^{pt}:\Re p<0\}$ and strictly proper rational functions with all poles in the open LHP.
\end{thmbox}
\begin{thmbox}[title={Theorem 0.29 (Lerch's uniqueness)}]
If $f,g$ satisfy (E) and $\Lap f(s)=\Lap g(s)$ for all $s>s_0$, then $f=g$ on $(0,\infty)$. \emph{This is what lets you invert by looking up the table.}
\end{thmbox}
\begin{methodbox}[title={Method: invert $F(s)$}]
\begin{enumerate}
\item If $F$ is not strictly proper, divide first (the polynomial part corresponds to $\delta$'s; in exam questions $F$ is usually strictly proper).
\item Factor the denominator. Write one term per factor: $\frac{A}{s-a}$ for a simple real root;
$\frac{A}{s-a}+\frac{B}{(s-a)^2}$ for a double root; $\frac{Bs+C}{s^2+bs+c}$ for an irreducible quadratic.
\item Find constants: \textbf{cover-up} (multiply by $(s-a)$, put $s=a$) for simple roots; compare coefficients or substitute convenient values for the rest.
\item Complete the square in quadratics: $s^2+2s+5=(s+1)^2+4$, and write the numerator in terms of $(s+1)$.
\item Read off from the table; cite Lerch.
\end{enumerate}
\end{methodbox}

\begin{workbox}[title={Worked example: repeated pole}]
$F(s)=\dfrac{s+3}{(s+1)^2(s+2)}=\dfrac{A}{s+1}+\dfrac{B}{(s+1)^2}+\dfrac{C}{s+2}$.
Cover-up at $s=-2$: $C=\frac{1}{1}=1$. Cover-up for $B$ (multiply by $(s+1)^2$, $s=-1$): $B=\frac{2}{1}=2$. Compare $s^2$ coefficients: $A+C=0\Rightarrow A=-1$.
So $f(t)=-e^{-t}+2te^{-t}+e^{-2t}$.
\end{workbox}

\subsection{Solving initial value problems \texorpdfstring{\fr{107--110}}{}}\label{sec:ivp}\label{ss:3-5}
\begin{plainbox}
Recipe: translate the whole ODE (initial values appear automatically), solve for $\hat y$ by simple algebra,
split into partial fractions, translate back. No guessing of particular integrals needed.
\par\smallskip\textbf{In the exam this looks like:} ``Solve the IVP by taking Laplace transforms'' (AE3, A3.2). Very likely in the exam.
\end{plainbox}
\begin{thmbox}[title={Proposition 0.30}]
For $y^{(n)}+a_{n-1}y^{(n-1)}+\dots+a_0y=u$ with initial values $y^{(j)}(0)=p_j$:
$(s^n+a_{n-1}s^{n-1}+\dots+a_0)\hat y(s)+q_{n-1}(s)=\hat u(s)$, with $q_{n-1}$ a polynomial of degree $\le n-1$ built from the $a_j$ and $p_k$.
The polynomial in front of $\hat y$ is the characteristic polynomial (of the companion matrix).
\end{thmbox}
\begin{methodbox}[title={Method: IVP by Laplace (5 steps)}]
\begin{enumerate}
\item Transform every term, using the derivative rule with the given initial values.
\item Collect $\hat y(s)$ on one side: $\hat y=\dfrac{\hat u(s)+(\text{initial-value polynomial})}{\text{characteristic polynomial}}$.
\item Partial fractions.
\item Invert with the table (cite uniqueness).
\item Check $y(0)$ (and $y'(0)$) and, if time, substitute into the ODE.
\end{enumerate}
\end{methodbox}

\begin{workbox}[title={Worked example (AE3, A3.2): $y'-7y=\sin2t$, $y(0)=0$}]
\textbf{Step 1.} $s\hat y-0-7\hat y=\dfrac{2}{s^2+4}$.\quad \textbf{Step 2.} $\hat y=\dfrac{2}{(s-7)(s^2+4)}$.\\
\textbf{Step 3.} $\dfrac{2}{(s-7)(s^2+4)}=\dfrac{A}{s-7}+\dfrac{Bs+C}{s^2+4}$. Cover-up: $A=\frac{2}{49+4}=\frac2{53}$.
Then $2=A(s^2+4)+(Bs+C)(s-7)$: $s^2$: $0=A+B\Rightarrow B=-\frac2{53}$; constant: $2=4A-7C\Rightarrow C=\frac{4A-2}{7}=\frac{8/53-106/53}{7}=-\frac{14}{53}$.\\
\textbf{Step 4.} $\hat y=\frac2{53}\cdot\frac1{s-7}-\frac2{53}\cdot\frac{s}{s^2+4}-\frac{7}{53}\cdot\frac{2}{s^2+4}$, so
\[\boxed{y(t)=\frac{2}{53}e^{7t}-\frac{2}{53}\cos2t-\frac{7}{53}\sin2t.}\]
\textbf{Step 5.} $y(0)=\frac2{53}-\frac2{53}=0$ \checkmark
\end{workbox}

\begin{workbox}[title={Worked example: second order, step input}]
$y''+3y'+2y=4$, $y(0)=1$, $y'(0)=0$.
\textbf{Step 1.} $(s^2\hat y-s)+3(s\hat y-1)+2\hat y=\frac4s$.
\textbf{Step 2.} $(s^2+3s+2)\hat y=\frac4s+s+3$, so $\hat y=\dfrac{s^2+3s+4}{s(s+1)(s+2)}$.
\textbf{Step 3.} Cover-up: at $s=0$: $\frac{4}{2}=2$; at $s=-1$: $\frac{1-3+4}{(-1)(1)}=-2$; at $s=-2$: $\frac{4-6+4}{(-2)(-1)}=1$.
\textbf{Step 4.} $y(t)=2-2e^{-t}+e^{-2t}$. \textbf{Check:} $y(0)=1$, $y'(0)=2-2=0$ \checkmark
\end{workbox}

\begin{workbox}[title={Worked example \fr{109--110}: transform only (third order)}]
$y'''+2y''-y'+y=u$, $y(0)=7$, $y'(0)=-1$, $y''(0)=5$:
$\Lap y'=s\hat y-7$, $\Lap y''=s^2\hat y-7s+1$, $\Lap y'''=s^3\hat y-7s^2+s-5$. Adding with coefficients $1,2,-1,1$:
$(s^3+2s^2-s+1)\hat y-7s^2-13s+4=\hat u(s)$.
\end{workbox}

\subsection{Convolution and integral equations \texorpdfstring{\fr{100--106}}{}}\label{sec:conv}\label{ss:3-6}
\begin{plainbox}
A \textbf{convolution} $\int_0^tf(t-u)g(u)\,du$ is a ``weighted memory'' of $g$. Its transform is simply the \textbf{product} $\hat f\,\hat g$.
So an integral equation with $\int_0^tk(t-u)y(u)\,du$ becomes an algebraic equation for $\hat y$.
\par\smallskip\textbf{In the exam this looks like:} ``Solve the integral equation'' or ``prove the convolution theorem''.
\end{plainbox}
\begin{thmbox}[title={Proposition 0.28 (Laplace convolution)}]
$(f*g)(x)=\displaystyle\int_0^xf(x-y)g(y)\,dy$ is commutative, linear, associative and
\[\boxed{\Lap(f*g)=\Lap(f)\,\Lap(g).}\]
\end{thmbox}
\begin{proofbox}[title={Proof of $\Lap(f*g)=\Lap f\,\Lap g$ \fr{101--102}}]
(E): if $|f|,|g|\le Me^{\beta x}$ then $|f(x-y)g(y)|\le M^2e^{\beta x}$, so $|f*g(x)|\le xM^2e^{\beta x}\le M^2e^{(\beta+1)x}$.
Swap the order of integration (region $0<y<x<\infty$) and substitute $u=x-y$:
\begin{align*}\int_0^\infty e^{-sx}\int_0^xf(x-y)g(y)\,dy\,dx&=\int_0^\infty e^{-sy}g(y)\Big(\int_y^\infty e^{-s(x-y)}f(x-y)\,dx\Big)dy\\&=\int_0^\infty e^{-sy}g(y)\,dy\int_0^\infty e^{-su}f(u)\,du.\end{align*}
\end{proofbox}
\begin{methodbox}[title={Method: integral equations}]
Spot the convolution: $\int_0^tk(t-u)y(u)\,du=(k*y)(t)$, which transforms to $\hat k(s)\hat y(s)$. Transform, solve the linear equation for $\hat y$, partial fractions, invert.
\end{methodbox}
\begin{workbox}[title={Worked example \fr{105--106}}]
$y(t)=2e^{at}+\int_0^te^{b(t-u)}y(u)\,du$ ($a\neq b+1$). Transform: $\hat y=\frac2{s-a}+\frac{\hat y}{s-b}$, so
$\hat y\,\frac{s-b-1}{s-b}=\frac{2}{s-a}$, i.e.\ $\hat y=\dfrac{2(s-b)}{(s-a)(s-b-1)}=\dfrac{A}{s-a}+\dfrac{B}{s-b-1}$ with
$A=\frac{2(a-b)}{a-b-1}$, $B=\frac{-2}{a-b-1}$ (cover-up). Hence $y(t)=\dfrac{2(a-b)e^{at}-2e^{(b+1)t}}{a-b-1}$.
\end{workbox}
\begin{workbox}[title={Worked example: $y(t)=t+\int_0^t\sin(t-u)\,y(u)\,du$}]
$\hat y=\frac1{s^2}+\frac{1}{s^2+1}\hat y\Rightarrow\hat y\cdot\frac{s^2}{s^2+1}=\frac1{s^2}\Rightarrow\hat y=\frac{s^2+1}{s^4}=\frac1{s^2}+\frac1{s^4}$, so $y(t)=t+\dfrac{t^3}{6}$.
\end{workbox}

\subsection{Laplace transform of \texorpdfstring{$(A,B,C,D)$}{(A,B,C,D)} \texorpdfstring{\fr{111--115}}{}}\label{sec:lapabcd}\label{ss:3-7}
\begin{plainbox}
Applying the translator to the whole system: integrating becomes dividing by $s$, so the matrix ODE becomes
$\hat X=(sI-A)^{-1}B\hat U$, and output $=T(s)\times$ input. This is \emph{why} $T(s)=D+C(sI-A)^{-1}B$.
\par\smallskip\textbf{In the exam this looks like:} ``Show that $\hat Y=T(s)\hat U$ when $X(0)=0$'' or ``find $\exp(tA)$ by inverting $(sI-A)^{-1}$''.
\end{plainbox}
\begin{thmbox}[title={Theorem 0.31}]
Transforming $\dot X=AX+BU$, $Y=CX+DU$ gives $s\hat X-X(0)=A\hat X+B\hat U$, $\hat Y=C\hat X+D\hat U$, so (for $s$ not an eigenvalue)
\[\hat X=(sI-A)^{-1}B\hat U+(sI-A)^{-1}X(0),\qquad \hat Y=\underbrace{\big(D+C(sI-A)^{-1}B\big)}_{T(s)}\hat U+C(sI-A)^{-1}X(0).\]
With $X(0)=0$: $\boxed{\hat Y(s)=T(s)\hat U(s)}$. The \emph{transfer function} of a system is the $T$ with $\hat Y=T\hat U$ (zero initial conditions).
\end{thmbox}
\begin{workbox}[title={Worked example: $\exp(tA)$ by Laplace, complex eigenvalues}]
$A=\begin{bmatrix}0&1\\-5&-2\end{bmatrix}$: $(sI-A)^{-1}=\dfrac{1}{s^2+2s+5}\begin{bmatrix}s+2&1\\-5&s\end{bmatrix}$ and $s^2+2s+5=(s+1)^2+4$.
\begin{itemize}
\item $\frac{s+2}{(s+1)^2+4}=\frac{s+1}{(s+1)^2+4}+\frac12\cdot\frac{2}{(s+1)^2+4}\ \to\ e^{-t}(\cos2t+\tfrac12\sin2t)$;
\item $\frac{1}{(s+1)^2+4}\to\frac12e^{-t}\sin2t$;\quad $\frac{-5}{(s+1)^2+4}\to-\frac52e^{-t}\sin2t$;
\item $\frac{s}{(s+1)^2+4}=\frac{s+1}{(s+1)^2+4}-\frac12\cdot\frac{2}{(s+1)^2+4}\to e^{-t}(\cos2t-\tfrac12\sin2t)$.
\end{itemize}
\[\exp(tA)=e^{-t}\begin{bmatrix}\cos2t+\frac12\sin2t&\frac12\sin2t\\-\frac52\sin2t&\cos2t-\frac12\sin2t\end{bmatrix}.\quad\text{Check }t=0:\ I\ \checkmark\]
\end{workbox}

\subsection{Jordan canonical form \texorpdfstring{\fr{116--126}}{}}\label{sec:jordan}\label{ss:3-8}
\begin{plainbox}
Not every matrix can be diagonalised. The \textbf{Jordan form} is the ``next best'': diagonal blocks with 1's just above the diagonal.
Each block with eigenvalue $\lambda$ gives terms $e^{\lambda t}, te^{\lambda t}, \dots$ in $\exp(tA)$. Count blocks = count eigenvectors.
\par\smallskip\textbf{In the exam this looks like:} ``Find the algebraic/geometric multiplicities and the minimal polynomial'' or ``write down $\exp(tA)$ for a Jordan block''.
\end{plainbox}
\begin{defbox}
Jordan block $J_k(\lambda)$: $\lambda$ on the diagonal, $1$ on the superdiagonal, zeros elsewhere.
\textbf{JCF (Theorem 0.32):} every $A$ is similar to a block-diagonal matrix of Jordan blocks: $A=S\,\mathrm{diag}(J_{k_1}(\lambda_1),\dots,J_{k_r}(\lambda_r))\,S^{-1}$.
\begin{itemize}
\item \textbf{Algebraic multiplicity} of $\lambda$ = power of $(z-\lambda)$ in $\chi_A$ = total size of the blocks with eigenvalue $\lambda$.
\item \textbf{Geometric multiplicity} = $\dim E(\lambda)=\dim\ker(A-\lambda I)$ = \textbf{number} of blocks with eigenvalue $\lambda$. So geometric $\le$ algebraic.
\item \textbf{Minimal polynomial}: $m(s)=\prod(s-\lambda)^{(\text{size of largest block for }\lambda)}$.
\item Generalised eigenvector: $v\neq0$ with $(\lambda I-A)^mv=0$ for some $m\ge2$.
\end{itemize}
\end{defbox}
\begin{workbox}[title={Worked example \fr{119--120} (corrected)}]
$A=J_3(2)\oplus J_1(2)\oplus J_2(3)$ (a $6\times6$ matrix: $3+1+2=6$).
$\chi_A(s)=(s-2)^4(s-3)^2$; $E(2)$ has dimension 2 (two blocks for $\lambda=2$), $E(3)$ dimension 1; eigenvectors $e_1,e_4,e_5$;
minimal polynomial $m(s)=(s-2)^3(s-3)^2$.\\
\textit{Note:} the lecture notes write the blocks as ``$J_2(3)\oplus J_2(1)\oplus J_2(3)$'' and ``eigenvalues $2,2,3$''; that is a typo. The matrix drawn, the characteristic polynomial and the minimal polynomial all match $J_3(2)\oplus J_1(2)\oplus J_2(3)$.
\end{workbox}
\begin{methodbox}[title={Method: find the Jordan form of a small matrix}]
\begin{enumerate}
\item Find $\chi_A$ and the algebraic multiplicities.
\item For each $\lambda$, find $\dim\ker(A-\lambda I)=n-\rank(A-\lambda I)$ = number of blocks.
\item If needed, $\rank(A-\lambda I)^2$ etc.\ tells you the block sizes. For $n\le3$ steps 1--2 usually decide everything.
\item $\exp(tJ_k(\lambda))=e^{\lambda t}\begin{bmatrix}1&t&\frac{t^2}{2!}&\cdots\\&1&t&\cdots\\&&\ddots&\\&&&1\end{bmatrix}$ (since $J_k(\lambda)=\lambda I+N_k$, $N_k^k=0$).
\end{enumerate}
\end{methodbox}

\subsection{Exam-style questions for Chapter 3}\label{ss:3-9}
\begin{answerbox}
\textbf{IVP:} (1) ``Taking Laplace transforms, using $\Lap(y')=s\hat y-y(0)$\dots''; (2) solve for $\hat y$; (3) partial fractions with constants found; (4) ``by the table and Lerch's uniqueness theorem, $y(t)=\dots$''; (5) check $y(0)$.\\
\textbf{From the definition:} write the integral to $R$, evaluate, let $R\to\infty$, state the condition on $s$.\\
\textbf{Integral equation:} name the convolution, transform, solve, invert, check by substitution.
\end{answerbox}
\begin{qbox}[title={Q3.1 (from the definition)}]
Show from the definition that $\Lap(e^{at})=\frac1{s-a}$ for $s>a$, and deduce $\Lap(\cosh at)$ and $\Lap(te^{at})$.
\end{qbox}
\begin{solbox}
$\int_0^Re^{at}e^{-st}dt=\Big[\frac{e^{-(s-a)t}}{-(s-a)}\Big]_0^R=\frac{1-e^{-(s-a)R}}{s-a}\to\frac1{s-a}$ since $s-a>0$.
$\cosh at=\frac12(e^{at}+e^{-at})\Rightarrow\frac12\big(\frac1{s-a}+\frac1{s+a}\big)=\frac{s}{s^2-a^2}$ ($s>|a|$).
$\Lap(te^{at})=-\frac{d}{ds}\frac1{s-a}=\frac1{(s-a)^2}$.
\end{solbox}

\begin{qbox}[title={Q3.2 (IVP)}]
Solve $y''+4y=8$, $y(0)=0$, $y'(0)=2$ by Laplace transforms.
\end{qbox}
\begin{solbox}
$s^2\hat y-2+4\hat y=\frac8s\Rightarrow\hat y=\frac{8}{s(s^2+4)}+\frac{2}{s^2+4}$. Partial fractions: $\frac{8}{s(s^2+4)}=\frac2s-\frac{2s}{s^2+4}$.
So $y=2-2\cos2t+\sin2t$. Check: $y(0)=0$ \checkmark; $y'=4\sin2t+2\cos2t$, so $y'(0)=2$ \checkmark; $y''+4y=8\cos2t-4\sin2t+8-8\cos2t+4\sin2t=8$ \checkmark
\end{solbox}

\begin{qbox}[title={Q3.3 (system with initial state)}]
For $A=\begin{bmatrix}0&1\\-2&-3\end{bmatrix}$, $B=\begin{bmatrix}0\\1\end{bmatrix}$, $C=[1\ 0]$, $D=0$, $u=0$, $X(0)=[1;0]$, find $\hat Y(s)$ and $y(t)$.
\end{qbox}
\begin{solbox}
$\hat Y=C(sI-A)^{-1}X(0)$; $(sI-A)^{-1}=\frac1{(s+1)(s+2)}\begin{bmatrix}s+3&1\\-2&s\end{bmatrix}$, so $\hat Y=\frac{s+3}{(s+1)(s+2)}=\frac{2}{s+1}-\frac1{s+2}$
and $y=2e^{-t}-e^{-2t}$ (matches the $(1,1)$ entry of $\exp(tA)$ in Chapter 2 \checkmark).
\end{solbox}

\begin{qbox}[title={Q3.4 (Jordan)}]
$A=\begin{bmatrix}3&1&0\\0&3&0\\0&0&3\end{bmatrix}$. Find the algebraic and geometric multiplicity of $3$, the minimal polynomial and $\exp(tA)$.
\end{qbox}
\begin{solbox}
$\chi_A=(s-3)^3$: algebraic multiplicity 3. $A-3I=\begin{bmatrix}0&1&0\\0&0&0\\0&0&0\end{bmatrix}$ has rank 1, so geometric multiplicity $3-1=2$: two blocks, $J_2(3)\oplus J_1(3)$.
Minimal polynomial $(s-3)^2$. $\exp(tA)=e^{3t}(I+tN)=e^{3t}\begin{bmatrix}1&t&0\\0&1&0\\0&0&1\end{bmatrix}$.
\end{solbox}
\newpage
